%!TEX root = main.tex

\section{Evaluated Dissemination Algorithms}
\label{sec:algorithms}

In this paper we focus on dissemination algorithms for wireless ad-hoc networks.
To ease the discussion about the algorithms that are evaluated hereafter, we rely on the taxonomy proposed  by Ruiz et al.~\cite{Ruiz2015}.
In that comprehensive survey, the authors present a large subset of dissemination algorithms and propose a taxonomy of features to classify them.
In addition, the survey makes a summary of the \textit{standard} terminology used in the literature.
The terms of interest for this paper follow: 

An ad-hoc network can be represented as a graph $G = (V, E)$, where $V$ is the set of
nodes composing it, and $E$ is the set of links connecting nodes in the communication range.


\begin{mydef}[Dissemination/Broadcast]
	Given a \textit{source node} $v$ starting a broadcast process, the broadcast problem consists in reaching any $u \in V$ with the message $m$ by following the set of links $e \in E$.
\end{mydef}

\begin{mydef}[Broadcast Session]
	In an ad-hoc network, nodes $V^\prime \subseteq V$ may aim at disseminating the messages $m_0, m_1, \dots, m_k$.
	The dissemination of a message $m_i$ is a \textit{broadcast session}; and it can be uniquely identified by the message.
\end{mydef}

\begin{mydef}[Underlying Topology]
	Dissemination algorithms often maintain a subset of nodes -- the underlying topology -- used to drive the broadcasting process. Those nodes, which play a special role during a broadcast session, are determined using strategies specifics to each dissemination technique.
\end{mydef}

In general, dissemination methods share a set of steps that include ``\textit{variation points}''.
The methods are then classified based on the concrete way in which these variation points are implemented.
For instance, the foremost variation point is the usage or not of an \textit{underlying topology} to drive the dissemination process.

Among algorithms that build an underlying topology, we study three in this paper:
multi-point relaying~\cite{Qayyum2002}, computing relay points based on connected dominating sets~\cite{Wu2001}, and a simplified version of a method that uses a tree-like topology~\cite{Cagalj2007}.
As stated, these methods, which are executed in every node $n \in V$, have a common structure that follows:

\begin{enumerate}
	\item \textbf{Compute underlying topology} $T$.
	\item When a broadcast message $m_i$ is received for the first time, process it.
	\item If $n \in T$ and \textbf{some other conditions} are met, \textbf{retransmit} $m_i$ to a subset of my direct neighbors.
\end{enumerate}

Notice how the variation points are highlighted in the previous steps.

On the other hand, methods that do not build an underlying topology can be classified using additional features.
For instance, how much contextual information they use.
In this case, we find \textit{context-aware} and \textit{context-oblivious} methods, but there are also methods where nodes only use a specific piece of contextual information (e.g., neighbor knowledge).
The algorithms of this group that we consider in this paper are context-aware and \textit{statistical}.
They share the following common structure, where $T$ is a value related to the context and $T_C$ is an associated threshold:

\begin{enumerate}
	\item \textbf{Initialize} $T$, \textbf{calculate threshold} value $T_C$.
	\item When a broadcast message $m_i$ is received, \textbf{update the value} of $T$.
		  
		  If $T < T_C$, do not retransmit and stop procedure. Otherwise, set random timer.
		  
	\item If $m_i$ is received again before the timer expires, repeat the previous step.
	\item When the timer expires, retransmit $m_i$ if $T > T_C$.
\end{enumerate}

\begin{figure*}[!h] 
	\centering
	\includegraphics[scale=0.3]{fig/taxonomy.png}
	\caption{Taxonomy to classify broadcasting algorithms}\label{fig:taxonomy}
\end{figure*}

In the rest of this section we briefly present the methods evaluated in the present work.
 
\subsection{Using underlying topology}

To build an underlying topology for a given method, the nodes in the network must exchange messages.
Some dissemination algorithms require relatively simple data from the context in order to build this topology.
However, in some cases the complexity of creating the underlying topology far exceeds that of disseminating messages.

A common type of information required is the set of neighbors of a node $v$ that can be reached using exactly a certain number $k$ of hops -- this is written as $N^{k}(v)$.
For instance $N(v)$ is the set of direct neighbors of $v$ (notice that no superscript is used when $k = 1$).
In the same way, $u \in N^2(v)$ if $\exists w \in V, \{(v, w), (w,u)\} \subseteq E$, but $(v, u) \notin E$.
This indicates the nodes reachable from $v$ using exactly one hop.

The \textbf{multipoint relaying  (MPR)} is an efficient technique for flooding in wireless networks~\cite{Qayyum2002}.
Using $N(v)$ and $N^2(v)$, each node determines if it is a member of the MPR set (it must retransmit broadcast messages).
To do so, it uses two rules:
\begin{enumerate}
	\item A node $w \in N(v)$ is member of MPR if $w$ is the only neighbor of a node in $N^2(v)$.
	\item While there exist a node $w \in N^2(v)$ which is not covered by a node in MPR:
		\begin{itemize}
			\item Select the node $z \in N(v)$, but not in MPR, that covers more uncovered nodes in $N^2(v)$. Add it to MPR.
		\end{itemize}
	\item Notify to each member of MPR that it must retransmit messages.
\end{enumerate}

Using this topology, a node simply retransmit a message the first time it is received.

Another technique determines the relay points using the \textbf{connected dominating sets (CDS)}~\cite{Wu2001}.
In this approach, a node retransmit a message if it is member of an approximation of the CDS.
Such an approximation is computed in a distributed manner by means of a proposed heuristic. 
Initially a node $v$ is marked as member of CDS if there exists two unconnected nodes $u, w \in N(v)$. This initial approximation can be improved using two rules

\newtheorem{myrule}{Rule}

\begin{myrule}
	Given nodes $v$, $u$ in $G^\prime$, 
	let $N\left[v\right] = N(v) \cup \{v\}$ and $N\left[u\right] = N(u) \cup \{u\}$,
	if $N\left[v\right] \subseteq N\left[u\right]$ and $id(v) < id(u)$
	then $G^\prime$ becomes $G^\prime - \{v\}$.
\end{myrule}

\begin{myrule}
	Assume that Given nodes $v$, $u$, $w$ in $G^\prime$. 
	if $N(v) \subseteq N(u) \cup N(w)$ and $id(v) = min \{ id(v), id(u), id(w) \}$, then $G^\prime$ is changed to $G^\prime - \{v\}$.
\end{myrule}

The above rules can be modified to take into account the battery level of each node.
In this way, the method can ensure that nodes will last longer.

\textbf{Embedded Wireless Multicast Advantage (EWMA)}~\cite{Cagalj2007} using~\cite{Gallager1983} to compute the minimum-weighted spanning tree.

\subsection{Without underlying topology}

The \textbf{distance-to-mean} (dist2mean) approach~\cite{Slavik2011} is based on a simple observation: if a node $v$ receives a message from all directions, there is little benefit in retransmitting it because likely all neighbors already received the message.
In other words, the neighborhood of $v$ is already covered.
To estimate the coverage of $v$, the method computes a metric for each broadcast session -- the distance to the spatial mean.
This is, every time the node $v$ receives a message, it records the position $(x_i, y_i)$ of the transmitting node.
When the timer expires and it is time to decide whether retransmit, this only happens if $T > T_C$, where:

\begin{equation}
T = \dfrac{\sqrt{(v_x-\bar{x})^2 + (v_y-\bar{y})^2}}{r}
\end{equation}

\begin{equation}\label{eq:dist2mean_threshold}
T_C = 0.95 -0.7e^{-0.11N}
\end{equation}

\begin{equation}
M = (\bar{x},\bar{y}) = (\dfrac{1}{n}\sum_{i=1}^{n}x_i, \dfrac{1}{n}\sum_{i=1}^{n}y_i)
\end{equation}

and $N$ is the number of neighbors of $v$ and $(v_x, v_y)$ is its location.
The value of function~\ref{eq:dist2mean_threshold}, which have been taken from~\cite{Slavik2011}, decreases as $N$ decreases in order to ensure connectivity.


The \textbf{area based beaconless (ABB)} algorithm~\cite{Ovalle-Martinez2006} uses a different approach to estimate whether a node $v$ must retransmit a message $m_i$.
As in the method \textit{dist2mean}, a metric is computed to determine if the transmission range of $v$ is already covered.
This metric, which is related to the context, is the angle covered by the transmitting nodes.
In practice, this can be expressed as the percentage $P$ of the communication perimeter that is covered by the received copies of $m_i$.
Another parameter of interest, the threshold $P_C$ is set to 100\%.
Finally, two methods to compute the timeout value of the timer are proposed: a uniform random value in the interval from 0 to 1; and the value $P/P_C$.
